\documentclass[11pt,letterpaper]{article}
\usepackage{fontspec,amsmath,amssymb,microtype}
\usepackage[margin=0.85in,headheight=17pt]{geometry}
\usepackage{graphicx,pdflscape,caption,xcolor,tikz}
\usetikzlibrary{arrows.meta,decorations.pathreplacing,positioning}
\usepackage{booktabs,tabularx,longtable,array,xurl,fancyhdr}
\usepackage[hidelinks,pdftitle={Linkwise Modem Implementation Reference}]{hyperref}
\definecolor{ovnavy}{HTML}{16324F}
\definecolor{ovblue}{HTML}{2E6F95}
\definecolor{ovteal}{HTML}{2A9D8F}
\definecolor{ovgold}{HTML}{E9C46A}
\definecolor{ovpale}{HTML}{EFF5F8}
\definecolor{ovgray}{HTML}{52606D}
\definecolor{ovred}{HTML}{B23A48}
\pagestyle{fancy}
\fancyhf{}
\lhead{\small LINKWISE \textbullet{} MODEM IMPLEMENTATION}
\rhead{\small Revision 0.1}
\cfoot{\thepage}
\setlength{\parindent}{0pt}
\setlength{\parskip}{5pt}
\setlength{\emergencystretch}{2em}
\newcommand{\modemcite}[1]{\cite{#1}\enspace}
\newcommand{\landscapediagram}[3]{%
  \clearpage\begin{landscape}\thispagestyle{plain}\centering
  #1\par\vspace{0.15in}
  \captionsetup{width=0.94\linewidth}
  \captionof{figure}{#2}\label{#3}
  \end{landscape}\clearpage
}
\newcommand{\landscapeplantuml}[3]{%
  \landscapediagram{\includegraphics[width=0.94\linewidth,height=0.76\textheight,keepaspectratio]{#1}}{#2}{#3}
}
\newcolumntype{Y}{>{\raggedright\arraybackslash}X}
\title{\textbf{Linkwise Modem Implementation Reference}\\[0.5em]
\large Waveforms, framing, and link recovery}
\author{Linkwise Project}
\date{Revision 0.1 \quad 12 November 2026}
\begin{document}
\maketitle
\begin{abstract}
This reference traces binary host input through frame assembly, 48 kHz
waveform generation, acquisition, decoding, and acknowledgment. It covers
BW500 and the wide index-modulated modes and works through a reproducible
BW500 exchange and a captured BW2300 recovery. Communications literature
supplies the theoretical models. Pinned Rust source, executable vectors,
and stock-peer recordings supply implementation evidence.
\end{abstract}
\clearpage
\tableofcontents
\clearpage

\section{Implementation boundaries}
The source baseline is \path{dcf1aa549f6affb6818e33ceb0ade51cd48dacb1}.
The wide sign-off used \texttt{70cf5ed}. The intervening commit records
that result in documentation. This reference shares the implemented DSP
treatment and audited BW500 exchange with the architecture white paper.
Measured protocol references and bench captures supply stock VARA
parameters \cite{stjohn,bench}.

\begin{table}[htbp]
\centering\small
\begin{tabularx}{\linewidth}{@{}p{0.20\linewidth}YY@{}}
\toprule
Boundary & Rust owner & Interface and state \\
\midrule
Host TCP & \path{api.rs}, \path{host/tcp.rs} &
Separate command and binary DATA streams. The DATA stream has no modem-frame boundaries. \\
Waveform session & \path{bin/linkwise-modemd/virtual_host.rs} &
Owns bench handshake, pending bytes, frame plan, turn, retries, and deadlines.
Audio records carry a big-endian sample count followed by little-endian float32 at 48 kHz. \\
Wire interpretation & \path{setup.rs}, \path{control.rs}, \path{data.rs} &
Packs caller identity, matches keyed controls, and validates full and closing DATA. \\
Signal processing & \path{receiver.rs}, \path{waveform.rs}, \path{bw500.rs},
\path{bw500_waveform.rs}, \path{wide_level4.rs}, \path{coding.rs} &
Acquires timing, produces soft decisions, and encodes and validates frames. \\
Validated-event runtime & \path{runtime.rs}, \path{arq.rs}, \path{audio/} &
Defines session, device presentation and release, and sender accounting contracts.\\
\bottomrule
\end{tabularx}
\caption{Implementation ownership from the host sockets to validated events.}
\end{table}

Host payload excludes protocol overhead. The wire body includes control
and any closing trailer and is followed by the CRC. The information block
may add pad bits before whitening. The transmitted codeword adds parity
and termination. Appendix~\ref{app:bw500-wire} gives the decoded byte
geometry and distinguishes encoder and decoder variants.

\section{BW500 digital signal processing}\label{sec:modem-dsp}
\input{modem-dsp-core}

\subsection{Max-log trellis metrics}
BCJR uses forward and backward recursions to combine transition evidence.
Turbo decoding exchanges soft information between two constituents.
Robertson, Villebrun, and Hoeher compare the max-log approximation
with log-MAP \cite{bcjr,berrou,robertson}. In \path{coding.rs}, let
\(u,p\in\{0,1\}\) be the input and parity of transition \(s\to s'\) at
index \(i\). The branch metric is
\[
 g_i(s,u)=(L_i^{\rm sys}+A_i)u+L_i^{\rm par}p .
\]
\(L_i^{\rm sys}\) and \(L_i^{\rm par}\) are unitless soft metrics, and
\(A_i\) is prior information from the other constituent. A bit of zero
adds no term. A bit of one adds its signed score. The recursions keep
the strongest compatible paths:
\begin{align}
 \alpha_{i+1}(s')&=\max_{s,u:\,s\to s'}(\alpha_i(s)+g_i(s,u)),\\
 \beta_i(s)&=\max_u(g_i(s,u)+\beta_{i+1}(s')),\\
 M_i(b)&=\max_{s,u=b}(\alpha_i(s)+g_i(s,u)+\beta_{i+1}(s')),\\
 E_i&=M_i(1)-M_i(0)-L_i^{\rm sys}-A_i.
\end{align}
The forward pass starts in state zero. The backward pass ends there.
Other boundary states start at negative infinity. The extrinsic score
compares the best paths for 1 and 0, then subtracts direct systematic
evidence and the prior so those terms are not counted again. If the
best scores are 8 and 3 and the systematic and prior scores are 2 and 1,
the illustrative extrinsic score is \(8-3-2-1=2\), favoring 1.
Appendix~\ref{app:bw500-wire} gives actual frame-1 demodulator values.

\section{Wide index-modulated waveforms}
\path{wide_level4.rs} implements BW2300 and BW2750 through a shared
\path{IndexWaveform} pipeline. Level 4 uses 92 bytes, \(N=736\) information
bits, and \(2N+12=1484\) coded bits. Level 3 uses 50 bytes plus two pad bits,
giving \(N=402\) and 816 coded bits. Level 2 uses \(N=200\) and
\(3N+12=612\) coded bits. Carrier allocations, pilot placement, and the
keyed lead come from measured protocol work \cite{stjohn}. Correlation
and tone hypotheses follow the signal-space framework in Proakis and
Salehi, chapters 3--4 \cite{proakis}.

\begin{table}[htbp]
\centering\small
\begin{tabular}{@{}rrrrrrll@{}}
\toprule
Level & Hop & Lead & Columns & Pilots & Bits/data col. & BW2300 bins & BW2750 bins \\
\midrule
2 & 1024 & 4 & 228 & 24 & 3 & 17--48 & 12--51 \\
3 & 1024 & 6 & 228 & 24 & 4 & 17--48 & 12--51 \\
4 & 512 & 12 & 395 & 24 & 4 & 9--24 & 6--25 \\
\bottomrule
\end{tabular}
\caption{Wide geometry at 48 kHz. Hop is in samples. Lead and column counts
are symbol periods. Columns include pilots.}
\end{table}

For allocation \(a_t\), Gray label \(G(\ell_t)\), stride \(d\), span \(S\),
and first bin \(k_0\), the DATA tone is
\[
 k_t=k_0+\bigl((a_t+dG(\ell_t))\bmod S\bigr),\qquad f_t=k_t\,48000/H .
\]
Select the base position, move by the Gray-mapped data offset, and wrap
within the available bins. \(H\) is samples per symbol and \(f_t\) is in
hertz. A BW2300 level-4 column with allocation and label both zero would
use bin 9, or 843.75 Hz. Lower modes use measured allocation arrays and
stride two. Level 4 uses generated allocations and stride one. The keyed
lead advances per emitted frame, including through the observed lower-rate
retry. The receiver searches a bounded window within two symbol lengths
of the energy candidate in 16-sample steps. Its timing score is the mean
of 24 normalized pilot powers. It rejects scores below 0.55.

For DATA bit position \(j\), let \(P_\ell\) be the correlation power of
candidate label \(\ell\). The wide decoder uses
\[
 m_j=\frac{8\bigl(\max_{\ell:\,\ell_j=1}P_\ell
               -\max_{\ell:\,\ell_j=0}P_\ell\bigr)}
             {\max(\epsilon,\max_\ell P_\ell)}.
\]
The numerator compares the best tone for each bit value. The denominator
scales the difference by the strongest tone in that column, with a floor
for an empty or very weak column. A positive value favors 1. For powers
of 9 and 3 in the two groups, the score is \(8(9-3)/9=5.33\).
The receiver scatters these scores into the common FEC order.
The setup frame uses level-4 geometry.

\section{Session control and recovery}
ARQ combines error detection with peer feedback and recovery \cite{lin}.
The virtual host matches keyed controls against its current role and
phase. The event-level sender consumes a validated answer for its
outstanding transmission token. A positive reply retires covered bytes.
A negative reply can retain the acknowledged prefix and repacketize the
outstanding range. Missing feedback triggers a state query in supported
contexts. A timeout can close the session with delivery unresolved.

The worked BW500 receiver accumulates DATA until a valid closing frame
completes the over. A repeated final query can elicit the stored closing
ACK without another host delivery. The stock rate-down replay replaces a
previously accepted 43-byte level-4 block with a level-3 block matching its
34-byte prefix. Sender \texttt{BUFFER 0}, receiver-delivered bytes, and a
Winlink mailbox record are distinct observations. \texttt{DISCONNECT}
can wait for queued DATA to drain. \texttt{ABORT} resets the session.
Protocol timeout can leave paired host sockets available for a fresh call.
Device presentation, turnaround, and release require operational-route
qualification.

\landscapeplantuml{build/diagrams/Linkwise_Modem_Recovery}{Event-level sender recovery. The session owner validates the wire response before advancing byte accounting. Appendix B gives the captured waveform-session sequence.}{fig:implementation-recovery}

\section{Verification and limits}
\begin{table}[htbp]
\centering\small
\begin{tabularx}{\linewidth}{@{}p{0.17\linewidth}p{0.24\linewidth}YY@{}}
\toprule
Concept & Literature & Rust implementation & Evidence \\
\midrule
Tone decisions & Proakis and Salehi, chs. 3--4 \cite{proakis}
& \path{waveform::Tone}, \path{bw500::scan}, \path{wide_level4}
& Published BW500 WAV replay and wide round trips. \\
Synchronization & Meyr et al., Part D \cite{meyr}
& \path{bw500_waveform::Waveform::decode}
& Measured-profile setup and PRBS9 DATA round trips. \\
CRC & Koopman and Chakravarty \cite{koopman}
& \path{coding::crc}, \path{data::decode}
& GENIBUS check vector and worked frame CRCs. \\
Iterative FEC & BCJR, turbo codes, max-log comparison \cite{bcjr,berrou,robertson}
& \path{coding::encode}, \path{coding::extrinsic}
& Corrupted systematic metrics recovered at both code rates. \\
ARQ & Lin et al. \cite{lin}, stock measurements \cite{stjohn}
& \path{arq::Sender}, \path{virtual_host}
& Withheld ACK, stock rejection, same-socket reconnect. \\
HF performance & F.1487, Annexes 1 and 3 \cite{itu}
& Controlled waveform fixtures and RF campaign plan
& Fixed echo tested. Calibrated fading sweeps pending. \\
\bottomrule
\end{tabularx}
\caption{Traceability from theory to Rust implementation and verification.}
\label{tab:modem-evidence-gates}
\end{table}

BW500 levels 1--3 recovered PRBS9, PRBS15, alternating bits, and constant
patterns through a fixed echo. Level 4 recovered the 64-byte PRBS9 demo
attachment across full and closing frames. The fixture uses the polynomial
\(x^9+x^7+x^6+x^4+1\) from MIL-STD-188-110C Appendix B \cite{mil}.
The manifest fixes register phase and byte packing. Winlink Express
exchanged both fixture files with stock VARA HF 4.9.1 in both caller roles
on virtual audio, with byte-exact and SHA-256 checks.

The wide DATA matrix met three consecutive live passes per tested
sender and receiver role at both bandwidths \cite{bench}. Remaining
gates include BW500 ACK-loss and mixed Pat repeats on the merged build,
native audio and PTT timing, and conducted and radiated tests. F.1487
supplies a Gaussian scatter model and test conditions for future fading
measurements, following the HF channel work of Watterson, Juroshek, and
Bensema \cite{watterson,itu}. In complex baseband, a tapped channel model
forms \(y(t)=\sum_i g_i(t)x(t-\tau_i)+n(t)\). Each path contributes a
delayed copy with a time-varying complex gain. The fixed-echo regression
uses a constant gain and one fixed delay, so it exercises a specific
receive case. Raw BER requires aligned decisions and declared
treatment of acquisition failures. Exact attachment delivery measures
the outcome after ARQ.

\appendix
\input{modem-wire-example}

\section{Worked BW2300 recovery against stock VARA}
Trial \path{tx-2300-02} used Linkwise as caller and stock VARA HF 4.9.1
on a Windows UTM guest \cite{bench}. The table uses the runner's relative
monotonic time. The Linkwise binary SHA-256 was
\path{f1218fa26fb7904265dd3953ad333972af6a94b48fefb4f1f0922d9ad058eb02}.

\begin{table}[htbp]
\centering\small
\begin{tabularx}{\linewidth}{@{}rp{0.27\linewidth}Y@{}}
\toprule
Time (s) & Event & Byte accounting \\
\midrule
23.823 & Delivery queued & 203 bytes, \texttt{00} through \texttt{ca}.\\
34.154 & First level-4 DATA & 89 bytes, \texttt{00}--\texttt{58}. Next \texttt{BUFFER} is 114.\\
34.486 & Guarded substitution & Next waveform replaced after its original SHA-256 matches.\\
46.490 & Linkwise state query & A stock reply precedes the lower-rate retry.\\
54.994 & Level-2 retry & 22 bytes, \texttt{59}--\texttt{6e}. 92 remain.\\
60.418 & Level-4 DATA & 89 bytes, \texttt{6f}--\texttt{c7}. Three remain.\\
65.852 & First delivery exact & Level-4 close carries \texttt{c8 c9 ca}. Total \(89+22+89+3=203\) bytes.\\
77.733 / 95.746 & Further deliveries & Exact six-byte reply and exact 203-byte reverse-ordered payload.\\
98.161 & First close & Both peers reported disconnect.\\
122.229 & Reconnect and close & Same sockets delivered an exact 12-byte probe. Both peers disconnected again.\\
\bottomrule
\end{tabularx}
\caption{Selected BW2300 recovery events. The pass covers 412 application
bytes plus the reconnect probe.}
\end{table}

Received role zero contains \texttt{F> 79} followed by carriage return.
Role one contains the ascending sequence, its reversal, and
\texttt{fresh link} followed by \texttt{00 ff}. The replacement PCM
SHA-256 is \path{5945af37ba584b8bf15de2db4569a09bdd4826077e48d425b9e1ec2eb6dff420}.
BW2750 sender runs used a stronger replacement and verified level-3
fallback. Both wide receiver rows observed a stock level-3 retry \cite{bench}.

\section{Reproduction}
Run from the checkout root. The profile helper validates pinned upstream inputs.
\begin{verbatim}
cargo test --locked --manifest-path rust/Cargo.toml --lib documented
cargo test --locked --manifest-path rust/Cargo.toml --lib wide
cargo test --locked --manifest-path rust/Cargo.toml --lib parity_repairs
profile=$(python3 tools/fetch-modem-profile.py)
LINKWISE_BW500_PROFILE="$profile" cargo test --locked \
  --manifest-path rust/Cargo.toml --test bw500_waveform -- --include-ignored
cargo build --locked --release --manifest-path rust/Cargo.toml \
  --bin linkwise-modemd
python3 tools/test-virtual-bw500.py rust/target/release/linkwise-modemd \
  "$profile" --wide
\end{verbatim}
With the local wide evidence set, run its saved \path{verify-evidence.py}
against that directory to check bytes and sign-off counts. A fresh live
trial requires the mixed bench procedure, stock VARA 4.9.1, and the
archived replacement stimulus.

\clearpage
\begin{thebibliography}{12}
\bibitem{proakis} J. G. Proakis and M. Salehi, \emph{Digital Communications},
5th ed., McGraw-Hill, 2008, chs. 3--5, 8, 11, and 13.
\href{https://glencoe.mheducation.com/sites/0072957166/information_center_view0/table_of_contents.html}{Publisher contents}.
\bibitem{meyr} H. Meyr, M. Moeneclaey, and S. A. Fechtel,
\emph{Digital Communication Receivers: Synchronization, Channel Estimation,
and Signal Processing}, Wiley, 1998, Part D.
\href{https://doi.org/10.1002/0471200573}{doi:10.1002/0471200573}.
\bibitem{bcjr} L. R. Bahl, J. Cocke, F. Jelinek, and J. Raviv,
Optimal decoding of linear codes for minimizing symbol error rate.
\emph{IEEE Trans. Inf. Theory}, 20(2), pp. 284--287, 1974.
\href{https://doi.org/10.1109/TIT.1974.1055186}{doi:10.1109/TIT.1974.1055186}.
\bibitem{berrou} C. Berrou, A. Glavieux, and P. Thitimajshima,
Near Shannon limit error-correcting coding and decoding: Turbo-codes.
\emph{Proc. IEEE ICC}, pp. 1064--1070, 1993.
\href{https://doi.org/10.1109/ICC.1993.397441}{doi:10.1109/ICC.1993.397441}.
\bibitem{robertson} P. Robertson, E. Villebrun, and P. Hoeher,
A comparison of optimal and sub-optimal MAP decoding algorithms operating
in the log domain. \emph{Proc. IEEE ICC}, pp. 1009--1013, 1995.
\href{https://elib.dlr.de/34203/}{DLR publication record}.
\bibitem{koopman} P. Koopman and T. Chakravarty,
Cyclic redundancy code (CRC) polynomial selection for embedded networks.
\emph{Proc. DSN}, 2004.
\href{https://users.ece.cmu.edu/~koopman/roses/dsn04/koopman04_crc_poly_embedded.pdf}{Author manuscript}.
\bibitem{lin} S. Lin, D. J. Costello Jr., and M. J. Miller,
Automatic-repeat-request error-control schemes.
\emph{IEEE Commun. Mag.}, 22(12), pp. 5--17, 1984.
\href{https://doi.org/10.1109/MCOM.1984.1091865}{doi:10.1109/MCOM.1984.1091865}.
\bibitem{itu} ITU-R, Recommendation F.1487-0,
\emph{Testing of HF modems with bandwidths of up to about 12 kHz using
ionospheric channel simulators}, 2000, Annexes 1 and 3.
\href{https://www.itu.int/rec/R-REC-F.1487-0-200005-I/en}{ITU publication}.
\bibitem{watterson} C. C. Watterson, J. R. Juroshek, and W. D. Bensema,
Experimental confirmation of an HF channel model.
\emph{IEEE Trans. Commun. Technol.}, 18(6), pp. 792--803, 1970.
\href{https://doi.org/10.1109/TCOM.1970.1090438}{doi:10.1109/TCOM.1970.1090438}.
\bibitem{mil} U.S. Department of Defense, MIL-STD-188-110C,
\emph{Interoperability and Performance Standards for Data Modems},
Appendix B, \S B.5.4.3.
\href{https://quicksearch.dla.mil/qsDocDetails.aspx?ident_number=35569}{DLA revision record}.
\bibitem{stjohn} S. St John (W9SSJ), \emph{hfmodem}, published VARA HF
protocol measurements, fixtures, and waveform constants, revision
\path{952f79245240132f4af209afbad1f4b2ba256613}.
\href{https://github.com/sstjohn/hfmodem/tree/952f79245240132f4af209afbad1f4b2ba256613/docs/protocols/vara}{Pinned measurement notes}.
\bibitem{bench} Linkwise Project, \emph{VARA HF interoperability bench},
November 12 wide DATA recovery sign-off.
\href{https://github.com/8r4n/linkwise/blob/dcf1aa549f6affb6818e33ceb0ade51cd48dacb1/docs/testing/vara-interoperability.md\#november-12-wide-data-recovery-sign-off}{Pinned sign-off record}.
\end{thebibliography}
\end{document}
